%% kapa: exam=1 round=24 year=2013 period=1 form=A subject=civil_law subject_ko=민법 seq=17 number=17
%% source: Q-Net 감정평가사 기출문제 (cst00309) artlSeq=5073802 file=extracted/1008/5073802/2013년 제24회 감정평가사 1차시험 기출문제/1교시(A형).hwp
%% conversion: kordoc 4.24.51 -> tools/gen1.py (auto); review=proofread-sonnet; flags=none
\begin{question}{17}
기한에 관한 설명으로 옳지 않은 것은? (다툼이 있으면 판례에 의함)
\begin{choices}
  \item 정지조건부 기한이익의 상실특약이 있는 경우 그 특약사유가 발생하더라도 채권자의 의사표시가 있어야 채무자는 기한의 이익을 상실한다.
  \item 파산선고를 받은 채무자는 기한의 이익을 주장하지 못한다.
  \item 채권자는 변제기까지의 이자를 포기하고 채무자에게 기한 전에 변제할 것을 청구할 수 없다.
  \item 종기있는 법률행위는 기한이 도래한 때로부터 그 효력을 잃는다.
  \item 당사자가 불확정한 사실이 발생한 때를 이행기한으로 정한 경우 그 사실의 발생이 불가능하게 된 때에도 이행기한이 도래한 것으로 보아야 한다.
\end{choices}
\end{question}
\answer{1}
